% Junior Level 3, Week 3 -- answer key
\answersweek{3}{Data Detectives}{Answers for workbook chapter 3 and play-at-home card 3.}

\section{Lesson questions}
\begin{answer}[Exit question]
Banana won, with 6. \enquote{banana} and \enquote{Banana} are the same answer written two ways, so 4 and 2 are added together: 6 is more than orange's 5.
\end{answer}

\begin{answer}[Reflection]
\enquote{Data is information we collect} (or \enquote{answers}, \enquote{things we count}) shows understanding. The tallest bar shows the answer most people gave. Any drawing of a fruit is right.
\end{answer}

\begin{answer}[Safety question]
Zoe should delete \enquote{What is your home address?}. An address is personal information, and the survey does not need it to find the favorite fruit.
\end{answer}

\section{Workbook exercises}
\begin{answer}[Exercise 3.1: Tally Time]
The tally shows apple 5, banana 3, orange 2, grapes 4. The bar graph has four bars that start at the bottom line: apple 5 squares tall, banana 3, orange 2, grapes 4. Apple won: circle the apple bar.

\textbf{What to look for:} \tickbox{} one square per tally mark \quad \tickbox{} no gaps at the bottom of a bar \quad \tickbox{} the child counts the gate of five as 5, not 4 or 6.
\end{answer}

\begin{answer}[Exercise 3.2: Fruit Survey Robot]
\textbf{What you should see:} after the green flag, both counters on the stage show 0. The cat asks \enquote{What is your favorite fruit?} and a typing box appears under the stage. Typing \enquote{apple} and pressing Return adds 1 to apple votes; any other answer adds 1 to other votes. After 4 answers the cat says \enquote{Thank you!}. The two counters always add up to 4.

\textbf{What to look for:} \tickbox{} \enquote{answer} sits in the left side of the \enquote{=} block and \enquote{apple} is typed on the right \quad \tickbox{} one \enquote{change by 1} in each gap of the if/else \quad \tickbox{} both counters are set to 0 at the start.
\end{answer}

\begin{answer}[Exercise 3.3: The Two Apples]
\textbf{The bug:} circle the rows \enquote{Apple} and \enquote{apple}. They are the same fruit, written once with a big A and once with a small a. Together they make 5 apple answers, so apple won, not banana (4). Fixed tally: apple 5, banana 4.

\textbf{On the Mac:} the robot counts \enquote{Apple} as apple: Scratch's \enquote{=} block does not care about big and small letters. It counts \enquote{apples} as other, because it is not exactly the same word. Any answer the child checked and wrote down is a good result; the point is that they tested it.
\end{answer}

\section{Play-at-home card}
\begin{answer}[Sock detectives]
Any tally and bar graph that match the socks are right: the bars together should have 10 squares, one for each sock. A story sentence like \enquote{Most of our socks are white} shows the child reads the graph.
\end{answer}

\begin{answer}[Good question, bad question]
Good: favorite animal, favorite color. No way: where you live and your password. Those are personal information, and a survey never needs them.
\end{answer}
